\chapter{مقدمه}
\newpage
در این فصل ابتدا تعریفی از مساله و سپس چالش‌های موجود در این حوزه مطرح می‌گردد. به دنبال چالش‌های مطرح شده در بخش‌های بعدی، اهداف دنبال شده در این رساله و جنبه‌ جدید بودن و نوآوری‌های رساله بیان می‌شود. در انتهای این فصل به ساختار رساله پرداخته شده است.   
‎\renewcommand{\baselinestretch}{2}\normalsize‎
\label{Introduction}
\section{تعریف مساله}
امروزه با افزایش حجم داده‌های پردازشی و محاسبات در بسیاری از کاربردها، چالش‌های بسیاری در فرآیند یادگیری الگوریتم‌های توزیع‌شده مطرح می‌شود. در یک شبکه تطبیقی توزیع‌شده\footnote{\lr{Distributed Adaptive Networks}} با محاسبات سنگین  یا داد‌های ذاتاً توزیع‌شده، موازی‌سازی فرآیند یادگیری و ایجاد یک مدل منسجم در کل شبکه اصلی‌ترین موضوع می‌باشد. به منظور تسریع عملیات و مدیریت راحت‌تر بخش‌های محاسبات و ارتباطات، یادگیری موازی باید تحت یک مدل محاسباتی موازی از جمله  \lr{PRAM}، \lr{LogP}، \lr{BSP} و \lr{MapReduce} پیاده‌سازی شود. هر الگوریتم موازی طراحی شده بر اساس مدل محاسباتی \lr{BSP} شامل تعدادی ابرگام با سه بخش مجزای محاسبه، ارتباط و همگام‌سازی مانعی می‌باشد. مقیاس‌پذیری، عملکرد بالا، استفاده آسان، پیش‌بینی زمان اجرا الگوریتم طراحی شده بر اساس آن و داشتن مدل برنامه‌نویسی از ویژگی‌های مهم این مدل محاسباتی است. مدل \lr{BSP} به راحتی قابل اعمال بر روی شبکه‌های تطبیقی توزیع‌شده است. این شبکه‌ها مقیاس‌پذیری بالایی دارند و در برابر خرابی‌ها مقاومت بوده و نیاز به سرور مرکزی ندارند. 

در یک محیط محاسباتی ناهمگن، گره‌های شبکه از لحاظ قدرت پردازشی و ارتباطی متفاوت خواهند بود و در صورت اعمال همگام‌سازی مانعی در مدل \lr{BSP}، مساله کندکننده‌ها رخ می‌دهد و میانگین زمان انتظار گره‌ها افزایش می‌یابد. در غیر این صورت، گره‌های کند می‌توانند شکاف تکرار را در بین گره‌های شبکه به تدریج افزایش داده و در نتیجه دقت یادگیری را کاهش دهند. به منظور همگرایی سریع و دقت یادگیری بالا، نیاز است که همگام‌سازی گره‌ها انجام شود اما باید سازگار با محیط‌های محاسباتی ناهمگن باشد.

با پیشرفت سریع فناوری‌های شبکه‌ای و ارتباطی، شبکه‌های توزیع‌شده‌ای که رایانه‌های شخصی، تلفن‌های همراه، حسگرها و محرک‌ها را به هم متصل می‌کنند، چهارچوب اصلی شبکه‌های ارتباطی و کنترل داده‌ها را تشکیل می‌دهند. در این شبکه‌ها نیاز به یادگیری توزیع‌شده زمانی به وجود می‌آید که یک هدف جهانی باید از طریق همکاری محلی بر روی تعدادی از عامل‌های به هم پیوسته محقق شود. این مساله در بسیاری از زمینه‌های مهم از جمله تخمین توزیع شده، یادگیری ماشین توزیع شده، تخصیص منابع و همچنین در مدل‌سازی رفتار دسته جمعی و ازدحامی توسط شبکه‌های بیولوژیکی رخ می‌دهد. شبکه‌های تطبیقی توزیع‌شده برای انجام وظایف پردازش و بهینه‌سازی اطلاعات غیرمتمرکز و مدل‌سازی انواع رفتارهای خودسازمان‌یافته و پیچیده‌ای که در طبیعت با آن مواجه هستیم، مناسب می‌باشند \cite{Zhao2014}. یک شبکه‌ تطبیقی توزیع‌شده شامل مجموعه‌ای از فیلترهای وفقی\footnote{\lr{Adaptive Filters}} یا عامل‌های یادگیری است. عامل‌ها از طریق توپولوژی اتصال به یکدیگر متصل می‌شوند و از طریق تعاملات محلی با یکدیگر همکاری می‌کنند تا مسائل بهینه سازی، تخمین و یا استنتاج توزیع‌شده را در زمان واقعی حل کنند. انتشار مداوم اطلاعات در سراسر شبکه توزیع‌شده، عامل‌ها را قادر می‌سازد تا عملکرد خود را در رابطه با جریان داده‌ها و شرایط شبکه تطبیق دهند. همچنین منجر به بهبود عملکرد سازگاری و یادگیری نسبت به عامل‌های غیرهمکار می‌شود. 

شبکه‌های تطبیقی توزیع‌شده جز دسته شبکه‌های غیرمتمرکز\footnote{\lr{Decentralized network}} می‌باشند و  با شبکه‌های متمرکز\footnote{\lr{Centralized network}} که یک گره مرکزی، به کلیه گره‌های شبکه متصل است و وظیفه تجمیع و پراکندگی اطلاعات محلی را بر عهده دارد، متفاوت می‌باشند. مقیاس‌پذیری و مقاومت در برابر خرابی گره‌ها و یال‌ها از ویژگی‌های مهم این شبکه‌ها است. همچنین اثربخشی و کارایی هر پیاده‌سازی تطبیقی توزیع‌شده در این شبکه‌ها به حالت همکاری بین عامل‌ها بستگی دارد. سه راه‌حل همکاری غیرمتمرکز مفید، از جمله استراتژی‌ اجماع\footnote{\lr{Consensus}}، استراتژی افزایشی\footnote{\lr{Incremental}} و استراتژی‌ انتشار\footnote{\lr{Diffusion}}، برای این منظور توسعه داده شده است \cite{Sayed2013} ، \cite{SayedBook}. روش همکاری انتشار نسبت به تغییرات محیط تطبیق‌پذیرتر از روش همکاری اجماع است و همانند روش همکاری افزایشی نیازی به یافتن دور در شبکه ندارد و هر گره می‌تواند با همه همسایگانش ارتباط داشته باشد و تخمین‌های دریافتی آنها را با داده‌های محلی خود تلفیق نماید. روش انتشار نسبت به اجماع دارای مزیت همگرایی سریع‌تر و رسیدن به میانگین مربع انحراف کمتر است \cite{Sayed2013,Sayed2014}. به طور خلاصه، روش‌های همکاری انتشار مقیاس‌پذیر، قوی، کاملاً توزیع‌شده هستند و به شبکه‌ها انطباق و توانایی یادگیری در زمان واقعی را می‌دهند. علاوه‌بر این، در مقاله \cite{Sayed2012} نشان داده شده است که روش‌های انتشار پایدارتر هستند و پایداری آن‌ها نسبت به توپولوژی شبکه غیر حساس است و منجر به بهبود کارایی و حالت پایدار در شبکه‌های تطبیقی می‌شوند.

الگوریتم‌های تطبیقی بسیاری وجود دارند که می‌توانند برای محاسبه و تخمین بردار پارامتر مسائل بهینه‌سازی مورد نظر در شبکه‌های تطبیقی توزیع‌شده مورد استفاده قرار بگیرند. برخی الگوریتم‌ها دقیق‌تر و گاهی پیچیده‌تر از سایرین هستند و برخی ساختاری ساده و در عین حال موثری دارند. بسیاری از این الگوریتم‌های تطبیقی مبتنی بر روش‌های گرادیان کاهشی\footnote{\lr{Gradient Descent}} می‌باشند. گرادیان کاهشی به روش‌های کمینه‌سازی اشاره دارد که محدود به استفاده از یک تابع هزینه\footnote{\lr{Cost function}}  خاصی نیستند. در حالیکه الگوریتم تطبیقی \lr{LMS} نوع خاصی از گرادیان کاهشی است که از تابع هزینه میانگین مربع خطا استفاده می‌کند. 

روش‌های گرادیان کاهشی\footnote{\lr{Gradient Descent}} الگوریتم‌ یادگیری مناسبی برای حل مسائل پیش‌بینی خطی بزرگ می‌باشند و یکی از محبوب‌ترین الگوریتم‌ها برای بهینه‌سازی و با اختلاف زیاد، پرکاربردترین روش برای بهینه‌سازی شبکه‌های عصبی است\cite{Bo2010}. روش‌های گرادیان کاهشی متفاوتی بر اساس حجم داده‌ استفاده شده برای محاسبه‌ی گرادیان تابع هدف \footnote{\lr{Cost function}} ارائه شده است. در روش گرادیان کاهشی تصادفی (\lr{SGD})\footnote{\lr{Stochastic Gradient Descent}} به ازای هر داده ورودی یک گام به‌روزرسانی بر روی بردار پارامتر انجام می‌شود. این روش تصادفی تغییرات را با واریانس بالایی ایجاد می‌کند که باعث می‌شود تابع هدف به شدت نوسان کند. برای حل این مشکل می‌شود به ازای هر دسته \footnote{\lr{Batch}} از داده‌های ورودی، یک گام گرادیان انجام شود. در این حالت واریانس تغییرات پارامترها کمتر می‌شود، در نتیجه نرخ همگرایی سریع‌تر خواهد بود. اندازه دسته می‌تواند برای کاربردهای مختلف متفاوت باشد\cite{We2018}.  

الگوریتم‌های گرادیان کاهشی تصادفی مرتبه دوم که جز این دسته از الگوریتم‌ها هستند، برای تخمین تطبیقی بردار پارامترهای مساله بهینه‌سازی در این رساله به‌کار گرفته می‌شوند. زیرا این الگوریتم‌ها برای هر دو نوع یادگیری برخط و برون‌خط کاربرد داشته و در حین اینکه حساس به ترتیب پردازش داد‌ه‌های ورودی نیستند، از دقت بالایی نیز برخوردار می‌باشند. با این اوصاف انتخاب مناسبی برای استفاده در شبکه‌های تطبیقی توزیع شده می‌باشند.

امروزه با افزایش حجم داده‌های پردازشی و نیاز به محاسبات سنگین در بسیاری از کاربردها، چالش‌های فراوانی در فرآیند یادگیری الگوریتم‌های تطبیقی مطرح شده است. در یک شبکه تطبیقی توزیع‌شده با حجم عملیات محاسباتی بالا یا حجم داده‌های پردازشی بالا، موازی‌سازی فرآیند آموزش و ایجاد یک مدل منسجم در کل شبکه  مهمترین مساله در فرآیند یادگیری خواهد بود. یادگیری موازی باید تحت یک مدل یا چارچوب محاسباتی موازی پیاده‌سازی و اجرا شود. از جمله این مدل‌های محاسباتی می‌توان به \lr{PRAM}، \lr{LogP}، \lr{BSP} و چارچوب محاسباتی \lr{MapReduce} اشاره کرد. مدل محاسباتی \lr{PRAM} (ماشین موازی با دستیابی تصادفی)\footnote{\lr{Parallel Random Access Machine}} \cite{Ja1992} یک مدل حافظه اشتراکی است که ارتباط بین پردازنده‌ها از طریق حافظه اشتراکی انجام می‌شود. به این دلیل که هزینه ارتباطات در این مدل بیشتر از هزینه محاسباتی است و تعداد پردازنده‌ها محدود است، این مدل برای پیاده‌سازی همه‌منظوره در سخت‌افزار مناسب نیست. چارچوب محاسباتی نگاشت-کاهش یا \lr{MapReduce} توسط گوگل\cite{De2008} توسعه یافته و برای حل برخی مسائل غیر پیش پا افتاده در زمینه‌های پردازش داده، داده‌کاوی و تجزیه و تحلیل گراف استفاده شده است. یک الگوریتم نگاشت-کاهش شامل تعدادی دور است که هر دور از یک فاز نگاشت و یک فاز کاهش تشکیل شده است. در این مدل ماشین‌ها تنها مجاز به برقراری ارتباط با ماشین پردازنده اصلی می‌باشند و حافظه اولیه محلی هر ماشین قبل از هر همگام‌سازی به طور کامل پاک می‌شود. در نتیجه برای پیاده‌سازی الگوریتم‌های تکرارشونده مناسب نمی‌باشد. مدل محاسباتی \lr{BSP} (موازی‌سازی همگام مانعی)\footnote{\lr{Bulk Synchronous Parallel}} \cite{Va1990} تعمیمی از مدل محاسباتی \lr{PRAM} است که توسط پروفسور \lr{L. G. Valiant } از دانشگاه هاروارد در سال 1990 پیشنهاد شده است. در سال‌های بعد، این مدل بوسیله پروفسور \lr{McCall} و \lr{Valiant} توسعه یافته است که به عنوان یک مدل پل برای محاسبات موازی برای سیستم‌هایی با حافظه توزیع‌شده کاربرد دارد. مدل محاسباتی \lr{BSP} شامل بخش‌های زیر می‌باشد:
\begin{itemize}
	\item
	مجموعه‌ای از کامپیوترها که هر یک دارای حافظه محلی می‌باشند
	\item 
	یک شبکه ارتباطی که پیام‌ها را از نقطه‌ای به نقطه دیگر انتقال می‌دهد
	\item 
	مکانیزمی که همگام‌سازی همه یا گروهی از پردازنده‌ها را انجام می‌دهد. 
\end{itemize}
برنامه‌های \lr{BSP} به صورت دنباله‌ای از ابرگام‌ها\footnote{\lr{Superstep}} نوشته می‌شوند. هر ابرگام از سه بخش اصلی تشکیل شده است: محاسبات بر روی داده‌های محلی، ارسال و دریافت پیام به منظور دسترسی به داده‌های غیرمحلی و همگام‌سازی مانعی به منظور تکمیل همه عملیات ارتباطی و از سرگیری محاسبات روی داده‌های محلی یا داده‌هایی که از طریق تبادل پیام به یک پردازنده رسیده‌اند. زمان اجرای یک رویکرد یا الگوریتم طراحی شده تحت مدل محاسباتی \lr{BSP} را می‌توان با استفاده از چهار پارامتر تاخیر شبکه، سرعت  پردازنده‌ها، تعداد پردازنده‌ها و نرخ راندمان ارتباطات پیش‌بینی کرد. 

این مدل محاسباتی منطبق با نیازها و ویژگی‌های شبکه‌های تطبیقی توزیع‌شده است و برای پیاده‌سازی الگوریتم‌های یادگیری تکرارشونده آن‌ها مناسب می‌باشد. مکانیزم همگام‌سازی مانعی در این مدل سبب می‌شود که در انتهای هر ابرگام تمام گره‌های شبکه عملیات ارتباطی و محاسباتی‌ خود را به اتمام رسانده و سپس  محاسبات ابرگام بعدی بر روی داده‌های محلی یا داده‌هایی که از طریق تبادل پیام به گره رسیده‌اند، انجام شوند. گره‌های شرکت‌کننده در یک شبکه تطبیقی توزیع‌شده ممکن است از لحاظ قدرت پردازشی و ارتباطی با یکدیگر متفاوت باشند. در نتیجه با یک محیط توزیع‌شده ناهمگن مواجه خواهیم بود که برخی از گره‌ها عملیات ارتباطی و یا محاسباتی را سریع‌تر انجام خواهند داد و به اجبار باید منتظر گره‌های کندتر در فرایند یادگیری بمانند تا پس از پایان عملیات آن گره‌ها بتوانند ابرگام بعدی را آغاز کنند. در چنین حالتی، مکانیزم همگام‌سازی مانعی سبب کاهش کارایی الگوریتم یادگیری و افزایش زمان انتظار گره‌های سریع‌تر در شبکه خواهد شد.

در این رساله یک رویکرد یادگیری توزیع‌شده برای تخمین پارامترهای مساله بهینه‌سازی در شبکه‌های تطبیقی توزیع‌شده ارائه می‌شود که مبتنی بر یک مدل محاسباتی توسعه‌یافته به نام \lr{BSP} فدرالی (\lr{FBSP}) است. همگام‌سازی گره‌های شبکه در این رویکرد با استفاده از یک تکنیک همگام‌ فدرالی انجام خواهد شد. در این تکنیک گره‌های شبکه در دو سطح منطقه‌ای و سراسری همگام می‌شوند. رویکرد پیشنهادی در دو مرحله اصلی برون‌خط و برخط اجرا می‌شود. مرحله برون‌خط که مرحله پیش‌پردازش نیز نامیده می‌شود، شامل دو بخش رصدکردن  شبکه و منطقه‌بندی آن است. در بخش رصد کردن شبکه، توان پردازشی هر گره محاسبه می‌شود. در بخش منطقه‌بندی شبکه، دسته‌بندی گره‌ها براساس فاصله گره‌ها از همدیگر، توان پردازشی و میزان اعتماد بین گره‌ها انجام می‌شود. در مرحله برخط، فرآیند یادگیری توزیع‌شده تحت مدل محاسباتی پیشنهادی \lr{FBSP} بر روی شبکه تطبیقی که در مرحله پیش‌پردازش منطقه‌بندی شده است، اجرا خواهد شد. مدل محاسباتی پیشنهادی همانند مدل \lr{BSP} شامل سه بخش مهم محاسبه، ارتباط و همگام‌سازی مانعی می‌باشد. با این تفاوت که مرحله ارتباط و همگام‌سازی مانعی در مدل \lr{FBSP} بر اساس یک روش منطقه‌بندی‌ انجام می‌شود که متناسب با ویژگی‌های شبکه‌ تطبیقی توزیع‌شده است. این روش منطقه‌بندی سبب می‌شود هزینه ارتباطات بین گره‌ها در زمان اجرا کاهش یابد و از وقوع مساله کندکننده‌ها جلوگیری به عمل می‌آید. در نتیجه زمان انتظار گره‌ها در شبکه کاهش خواهد یافت.

مساله کندکننده‌ها به یکی از دو دلیل 1) وجود گره‌های کند در بین گره‌های شبکه و 2) عدم توازن بار در بین گره‌‌ها، رخ خواهد داد. در حالتی که داد‌ه‌های پردازشی در شبکه توزیع‌شده از نوع ویدیو یا هر داده دیگری با اندازه‌های متفاوت باشد، عدم توازن بار را در گره‌ها خواهیم داشت که سبب می‌شود برخی از گره‌ها با اندازه داده یا اندازه ویدیو کوچک‌تر محاسبات‌شان زودتر تمام شود و منتظر سایر گره‌ها برای اتمام کارشان بمانند. در صورتی‌که توزیع بار بین گره‌ها به صورت یکسان انجام شده باشد، با افزایش تعداد گره‌هاي شرکت‌کننده در محیط توزیع‌شده، احتمال حضور گره‌های کند و گره‌هایی با قدرت پردازشی  و قدرت ارتباطی متفاوت در شبکه بسیار بالا خواهد رفت و یک محیط ناهمگن داریم. این موضوع افزایش زمان انتظار گره‌های سریع‌تر را نتیجه خواهد داد. در نتیجه زمان انتظار و زمان بیکار ماندن گره‌ها افزایش پیدا می‌کند.

رویکرد یادگیری توزیع‌شده پیشنهادی که یک رویکرد غیرمتمرکز است در واقع یادگیری مشارکتی را پیاده‌سازی می‌نماید. راه‌حل‌های غیرمتمرکز پتانسیل بالایی برای نسل پنجم (\lr{5G}) صنعتی از جمله رانندگی مشارکتی در وسایل نقلیه خودکار متصل و روباتیک مشارکتی در تولید هوشمند دارند. رویکرد پیشنهادی برای دو کاربرد شناسایی سیستم و پیش بینی بازار بورس مورد بررسی قرار گرفته است. کاربرد این رویکرد به این دو مساله محدود نمی شود بلکه قابل اعمال بر روی کلیه مسائل یادگیری با هدف تخمین پارامترها و مسائل پیش‌بینی مختلف خواهد بود. در مساله شناسایی سیستم هدف شناسایی مدل ریاضی سیستم ناشناخته با استفاده از داده‌های اندازه‌گیری‌شده است. در بازار بورس روزانه با حجم زیادی از اطلاعات و تراکنش‌ها مواجهه خواهیم بود و گاها نیاز به تحلیل‌های آنی و سریع وجود دارد تا‌‌ با پیش‌بینی‌های دقیق امکان انجام معامله‌های هوشمندانه و پرسود میسر شود. مساله پیش‌بینی بازار بورس سهام از جمله مسائلی است که رویکرد یادگیری توزیع‌شده پیشنهادی قادر است پیش‌بینی را با دقت بالا و نرخ خطا پایین انجام دهد.  

\section{چالشهای موجود}
در این بخش به بررسی چالش‌های موجود در مساله که به موجب آنها سمت و سوی اصلی این رساله شکل گرفته است، پرداخته شده است:
\begin{itemize}
	\item
	%1
	با رشد سریع حجم داده‌ها و محاسبات، فرآیند یادگیری در شبکه‌های تطبیقی توزیع‌شده نیازمند یک ساختار استاندارد موازی با تسریع هدفمند است. براساس بررسی‌های انجام‌شده در بخش قبل، یکی از چالش‌های اساسی در شبکه‌های تطبیقی توزیع‌شده عدم ارائه یک رویکرد یادگیری توزیع‌شده مبتنی بر یک مدل محاسباتی موازی برای این نوع از شبکه است که قادر باشد تعاملات بین گره‌ها را به خوبی در دنیای واقعی مدیریت نماید و متناسب با خصوصیات و نحوه ارتباطات و تعاملات گره‌ها در این نوع شبکه‌ها باشد و قابلیت پیش‌بینی زمان اجرا را در زمان پیاده‌سازی رویکرد داشته باشد.
	
	‌\item 
	%2
	هنگامی‌که پیاده‌سازی شبکه‌های تطبیقی توزیع‌شده به یک محیط بزرگ و ناهمگن گسترش می‌یابد، ناهمگنی قطعی در سخت‌افزار و قدرت محاسباتی گره‌ها را دیگر نمی‌توان نادیده گرفت. همچنین، در محاسبات ناهمگن توزيع‌شده مانند رایانش ابری، سرعت پردازش گره‌ها از لحاظ محاسباتی يا ارتباطی اغلب متفاوت است، از اين رو الگوريتم‌هاي یادگیری همگام در اين محيط‌ها سبب كاهش نرخ همگرايي و در بعضی مواقع واگرايي مي‌شوند. یکی دیگر از چالش‌های موجود، ارائه یک رویکرد یادگیری توزیع‌شده سازگار با محیط‌های ناهمگن و بارهای ‌کاری نامتوازن در شبکه است. 
	
	\item 
	%3
	در روش‌های مختلف یادگیری توزیع‌شده، همگام‌سازی گره‌ها بخش اصلی فرآیند یادگیری را تشکیل می‌دهد و تاثیر به‌سزایی بر روی مقیاس‌پذیری بهتر در محیط‌های توزیع‌شده ناهمگن دارد. در نتیجه، نیاز به یک تکنیک همگام‌سازی مدل‌ها در چنین محیط‌هایی به شدت احساس می‌شود که هم نرخ‌ همگرایی مناسبی داشته باشد و هم هزینه ارتباطات آن قابل قبول باشد. این شبکه‌ها که غيرمتمركز بوده و فاقد سرور پارامتري براي همگام‌سازي و تجميع مدل‌هاي برآورده شده هستند، نیازمند روشي براي همگام‌سازي با توجه به مساله و كاربرد می‌باشند. به طور معمول، همگام‌سازی مدل‌ها در گره‌هاي شبكه در انتهای هر تكرار انجام می‌شود. البته اجرای زياد همگام‌سازی باعث افزایش زمان یادگیری و کاهش بازدهی می‌گردد.

	‌\item 
	%4
	از دیگر چالش‌های موجود در رویکردهای توزیع‌شده، برقراری مصالحه بین پیچیدگی ارتباطی و نرخ همگرایی است. در یادگیری توزیع‌شده محاسبات در چند گره اجرا می‌شوند ولی سرعت اجرا الزاماً برابر با تعداد گره‌ها ضربدر سرعت یک گره نیست و این به دلیل هزینه ارتباطات بین گره‌ها است. در نتیجه، در طراحی تکنیک‌های همگام‌سازی توجه به میزان و پیچیدگی ارتباطات و هزینه‌های آن در مقابل حفظ همگرایی از الزامات کار می‌باشد.
	
		\item 
		%5
	‌یکی از چالش‌‌ها عدم وجود یک الگوریتم انتشاری تطبیقی همه منظوره برای شبکه‌های تطبیقی توزیع‌شده است که علاوه بر مقاوم بودن در برابر انواع نویزهای محیطی و محیط‌های ایستا و غیرایستا، قابلیت این را نیز داشته باشد که در محیط‌های برون‌خط و برخط قابل اجرا باشد و نسبت به ترتیب ورود داده‌ها حساس نباشد.
\end{itemize}
در این رساله سعی  شده است با مطالعات دقیق و بررسی‌های همه‌جانبه بر روی خصوصیات شبکه‌های تطبیقی توزیع‌شده، یک رویکرد یادگیری توزیع‌شده برای این شبکه‌ها بر پایه یک مدل محاسباتی توسعه‌یافته ارائه شود که گره‌ها به صورت موازی اقدام به یادگیری توزیع‌شده انتشاری نموده و بر اساس تکنیک همگام‌سازی مناسبی اقدام به‌روزرسانی مدل خویش نمایند.

\section{اهداف مساله}
با توجه به چالش‌های بیان شده در بخش قبل، در این رساله سعی شده است رویکردی ارائه شود تا فرآیند یادگیری و تخمین پارامتر را در یک شبکه تطبیقی توزیع‌شده‌ ناهمگن تحت یک مدل محاسباتی موازی توسعه‌یافته انجام دهد، به گونه‌ای که اهداف زیر مقرر گردد:
\begin{itemize}
	\item
	%1
	با توجه به چالش‌های موجود، ارائه یک مدل محاسباتی موازی جهت اجرا در شبکه‌های تطبیقی توزیع‌شده‌ به شدت مورد نیاز است. در این رساله، یکی از اهداف اصلی ارائه مدلی محاسباتی موازی متناسب با ویژگی‌های این شبکه‌ها برای حل مسائل بهینه‌سازی در یک محیط ناهمگن است. ناهمگنی در یک محیط محاسباتی به دلایل تفاوت در قدرت پردازشی و ارتباطی گره‌های شبکه یا  بارهای کاری نامتوازن رخ می‌دهد.	
	\item 
	%2
	با توجه به اینکه در یک محیط اجرایی ناهمگن ممکن است گره‌های شبکه از لحاظ قدرت پردازشی و ارتباطی متفاوت باشند، یکی از اهداف مورد نظر در این رساله استفاده از یک روش خوشه‌بندی محتوایی-ساختاری برای منطقه‌بندی گره‌های شبکه است. در این راستا، منطقه‌بندی گره‌ها در شبکه بر اساس یک ماتریس شباهت انجام می‌شود که با استفاده از ویژگی‌های ساختاری و محتوایی گره‌ها مقداردهی می‌شود. در این رساله، ویژگی‌های ساختار هر گره شامل فاصله آن گره از سایر گره‌ها و ضریب اعتماد به داده‌های دریافتی از گره‌های دیگر است و ویژگی‌های محتوایی هر گره شامل قدرت پردازشی و ارتباطی آن گره می‌باشد.  

	\item 
	%3
	با توجه به اینکه یک شبکه تطبیقی توزیع‌شده در واقع یک شبکه غیرمتمرکز و نظیر-به-نظیر با قابلیت تعامل و همکاری بین گره‌ها است، همگام‌سازی مدل‌های تخمین‌زده شده توسط گره‌ها امری بسیار مهم و موثر در کارایی و همگرایی مدل‌ها است و ناهمگنی محیط اجرا می‌تواند سبب افزایش زمان انتظار گره‌ها در شبکه در حین فرآیند یادگیری و کاهش دقت تخمین و گاها سبب عدم همگرایی شود. به همین منظور، لازم است در هر مرحله یادگیری، همگام‌سازی مانعی در محیط اجرایی ناهمگن در یکی از دو سطح منطقه‌ای یا سراسری انجام می‌شود. انجام همگام‌سازی مانعی دو سطحی علاوه بر افزایش نرخ همگرایی، سبب مقابله با مساله کندکننده‌ها و کاهش زمان انتظار گره‌ها می‌شود. 
	\item 
	%4
	یکی از چالش‌ها برقراری مصالحه بین پیچیدگی ارتباطی و نرخ همگرایی است. یکی از اهداف این رساله کاهش پیچیدگی و میزان ارتباطات گره‌ها در رویکرد یادگیری توزیع‌شده است. بدین منظور، به دنبال این هستیم که ارتباط گره‌ها بر اساس منطقه‌بندی گره‌ها در دو سطح منطقه‌ای و سراسری انجام شود. گره‌هایی که بر اساس معیارهای خوشه‌بندی در یک منطقه مشابه قرار می‌گیرند، تنها با گره‌های همسایه در آن منطقه ارتباط برقرار می‌کنند و مدل‌ پارامتری تخمین‌زده خویش را با گره‌ها تبادل خواهند کرد. اما در سطح سراسری، گره‌ها با همسایه خود در کل شبکه به ارتباط و تبادل اطلاعات می‌پردازند. منطقه‌بندی گره‌ها و انجام ارتباطات در دو سطح سبب کاهش هزینه‌های ارتباطی نیز خواهد شد. 
	\item 
	%5
	حضور انواع مختلف نویز در محیط‌های توزیع‌شده غیر قابل اجتناب است. یکی از اهداف مورد نظر در این رساله، ارائه الگوریتم یادگیری توزیع‌شده‌ای است که در مقابل انواع نویزهای گاوسی و غیرگاوسی مقاوم باشد. همچنین کارایی و همگرایی الگوریتم ارائه شده هم در محیط‌های ایستا و هم غیرایستا قابل قبول باشد. بدین منظور یک الگوریتم یادگیری توزیع‌شده انتشاری ارائه خواهد شد که در انواع محیط‌ها و در برابر انواع نویزها مقاوم خواهد بود. 
	\item 
	%6
	یکی از اهداف دیگر که در این رساله به دنبال آن هستیم، ارائه ابزاری یکپارچه با استفاده آسان برای برنامه‌های توزیع‌شده در بستر شبکه‌های تطبیقی است که مستقل از سخت‌افزار و سریع بوده و قابل اجرا بر روی کلیه ماشین‌ها باشد. 
	
\end{itemize}


\section{جنبه جدید بودن و نوآوری}

جنبه‌های جدید بودن و نوآوری راهکار پیشنهادی به شرح زیر می‌باشد:

\begin{itemize}
	\item
	%1
	رویکرد یادگیری توزیع‌شده دو مرحله‌ای: یک شبکه‌ تطبیقی توزیع‌شده در واقع یک گراف نظیر-به-نظیر است که تمام گره‌ها در سطح برابری از مسئولیت قرار دارند و هیچ گره مرکزی به منظور کنترل و مدیریت فعالیت سایر گره‌ها در این نوع شبکه‌ها وجود ندارد. در نتیجه در هنگام ارائه یک رویکرد یادگیری توزیع‌شده برای این شبکه‌ها توجه به این مساله بسیار مهم است. به همین دلیل در زمان اجرا به ندرت می‌توان تصمیم‌گیری پویا و سراسری انجام داد. در این رساله یک رویکرد یادگیری توزیع‌شده برای این نوع شبکه‌ها در دو مرحله برون‌خط و برخط معرفی شده است. در مرحله برون‌خط رصد کردن شبکه و عملیات پیش‌پردازشی موردنیاز انجام می‌شود و در مرحله برخط فرایند یادگیری تکرارشونده به صورت توزیع‌شده بر روی گره‌ها اجرا می‌شود.
	\item 
	%2
	ارائه یک الگوریتم یادگیری سازگار با محیط‌های ناهمگن: هنگامی‌که پیاده‌سازی شبکه‌های توزیع‌شده به یک محیط بزرگ گسترش می‌یابد، قطعا ناهمگنی در سخت‌افزار و قدرت محاسباتی گره‌ها را نمی‌توان نادیده گرفت. کارهای پیشین عمدتاً در یک محیط اجرای همگن مؤثر هستند، اما برای محیط‌های پردازشی که شامل گره‌هایی با قدرت محاسباتی و ارتباطی متفاوت هستند، مناسب نمی‌باشند. در این رساله سعی شده است تا الگوریتمی سازگار با این محیط‌ها ارائه شود.
	\item
	%3
	مدل محاسباتی موازی توسعه‌یافته برای محیط‌های ناهمگن: در تمامی مطالعات پیشین، نیاز به یک مدل محاسباتی موازی برای شبکه‌های تطبیقی توزیع‌شده و مناسب محیط‌های اجرایی ناهمگن احساس می‌شود. در این رساله سعی شده است یک مدل محاسباتی موازی توسعه‌یافته به نام \lr{FBSP} منطبق با ویژگی‌های شبکه‌های تطبیقی توزیع‌شده در محیط‌های ناهمگن معرفی گردد. در این مدل محاسباتی فازهای ارتباطات و همگام‌سازی مانعی بر اساس یک روش منطقه‌بندی‌ انجام می‌شوند تا هزینه ارتباطات و احتمال رخ‌دادن مشکل کندکننده‌ها کاهش یابد.
	\item
	%4
	منطقه‌بندی گره‌های شبکه: جهت داشتن یک همگام‌سازی غیرهم‌زمان با هدف حفظ نرخ همگرایی و کاهش پیچیدگی ارتباطی، منطقه‌بندی گره‌های شبکه با استفاده از یک روش خوشه‌بندی افزایشی و با توجه به نرخ محاسبه، همسایگی و میزان اعتماد آنها به همدیگر انجام می‌شود. 
	\item
	%5
	توسعه یک الگوریتم یادگیری تطبیقی انتشاری مبتنی بر کاهش گرادیان تصادفی: این دسته از الگوریتم‌های گرادیان کاهشی هم برای یادگیری برخط و هم یادگیری برون‌خط کاربرد داشته و حساسیت کمی به ترتیب پردازش داد‌ه‌های ورودی دارند، در این پژوهش سعی شده است این دسته از الگوریتم‌ها را برای برآورد بردار پارامترهای مجهول مساله بهینه‌سازی توسعه داد و الگوریتم قوی و مقاومی در برابر انواع نویزهای ورودی و انواع محیط‌های ایستا و غیرایستا برای یادگیری توزیع‌شده انتشاری معرفی کرد. 
	\item 
	%6
	ارائه یک ابزار یکپارچه نرم‌افزاری برای توسعه برنامه‌های توزیع‌شده تطبیقی: در این رساله سعی شده است به منظور استفاده عموم یک ابزار نرم‌افزاری به نام \lr{FBSPonMPI} ارائه شود که ضمن امکان استفاده آسان و سریع بودن، مستقل از سخت‌افزار باشد و بر روی همه ماشین‌های دارای \lr{MPI} قابل اجرا باشد. 
\end{itemize}


\section{ساختار رساله}
این رساله شامل شش فصل است. در فصل اول، کلیات، اهداف و ایده‌های کلی برای مقابله با چالش‌ها بیان گردید. در فصل دوم، ابتدا یک نگاه اجمالی به مفاهیم اولیه شبکه‌های تطبیقی توزیع‌شده خواهیم داشت و پس از آن مدل‌های محاسباتی موازی موجود برای پردازش موازی و توزیع‌شده داده‌ها را مورد بررسی و مقایسه قرار می‌دهیم. در فصل سوم، کارهای انجام شده از نقطه نظرهای مختلف بررسی  می‌شوند. در فصل چهارم، راهکار پیشنهادی را به صورت کامل تشریح خواهیم کرد و در فصل پنجم به ارزیابی و تحلیل نتایج بدست آمده و مقایسه راهکار ارائه شده با دیگر روش‌ها پرداخته می‌شود. در نهایت در فصل ششم، نتیجه‌گیری و مسیرهای اصلی کارهای آینده بیان خواهد شد.
